\begin{abstract}
We study the design of airline hub-and-spoke networks through a continuous optimization framework that jointly models network deployment, passenger choice, and market profitability. The lower-level model combines operating costs, sparse deployment penalties, price-weighted passenger utility terms, and entropy regularizers that recover a logit demand split between the new airline and the outside alternative. This yields a continuous approximation of hub activation and route opening decisions while preserving the economic structure of the airline design problem.

On top of this model, we introduce a bilevel formulation in which market weights are learned rather than fixed. The upper level maximizes the real objective of the airline, namely profit, by adjusting market-specific hyperparameters that scale the lower-level utility and entropy terms. This allows the model to emphasize the origin-destination markets that are most relevant for profitability, instead of relying on exogenous weights chosen \emph{a priori}.

The main difficulty is the simultaneous presence of hierarchical optimization and fixed deployment costs. To address it, we combine a reweighted $\ell_1$ approximation of fixed costs with a penalty-based bilevel strategy. The resulting algorithm alternates an outer sparsity loop, which updates the reweighting coefficients associated with capacities, and an inner penalty loop, which alternates between a penalized bilevel solve and a lower-level solve to update the market hyperparameters. The final methodology provides a practical route to solving airline network design problems that are beyond the reach of exact mixed-integer bilevel formulations, while maintaining an explicit link between passenger behavior, network topology, and profit.
\end{abstract}
